\documentclass[10pt,a4paper]{amsart}
\usepackage{amsmath,amssymb,amsthm,graphicx}
% Retained minimal preamble for the Stacks Project chapter more-algebra.tex (More on Algebra).
% The chapter file pinned-original-source.tex carries only \input{preamble} on its first line, so no
% preamble text lives inside the chapter itself.  The Stacks Project uses one single shared preamble for
% every chapter; the snapshot stages it identically as sources/stacks-chow/preamble.tex and
% sources/stacks-etale-cohomology/preamble.tex.  The block below is copied mechanically, line for line,
% from sources/stacks-chow/preamble.tex lines 14-18 (the Xy-pic setup) and 154-217 (the native theorem
% declarations together with the mathematical macros).  Nothing was added, removed or redefined; the
% wrapper class and the amsmath/amssymb/amsthm/graphicx base packages are supplied by the assembler.
% Chapter: more-algebra.tex -- More on Algebra.
% For commutative diagrams we use Xy-pic
\usepackage[all]{xy}

% We use 2cell for 2-commutative diagrams.
\xyoption{2cell}

% Theorem environments.
%
\theoremstyle{plain}
\newtheorem{theorem}[subsection]{Theorem}
\newtheorem{proposition}[subsection]{Proposition}
\newtheorem{lemma}[subsection]{Lemma}

\theoremstyle{definition}
\newtheorem{definition}[subsection]{Definition}
\newtheorem{example}[subsection]{Example}
\newtheorem{exercise}[subsection]{Exercise}
\newtheorem{situation}[subsection]{Situation}

\theoremstyle{remark}
\newtheorem{remark}[subsection]{Remark}
\newtheorem{remarks}[subsection]{Remarks}

\numberwithin{equation}{subsection}

% Macros
%
\def\lim{\mathop{\mathrm{lim}}\nolimits}
\def\colim{\mathop{\mathrm{colim}}\nolimits}
\def\Spec{\mathop{\mathrm{Spec}}}
\def\Hom{\mathop{\mathrm{Hom}}\nolimits}
\def\Ext{\mathop{\mathrm{Ext}}\nolimits}
\def\SheafHom{\mathop{\mathcal{H}\!\mathit{om}}\nolimits}
\def\SheafExt{\mathop{\mathcal{E}\!\mathit{xt}}\nolimits}
\def\Sch{\mathit{Sch}}
\def\Mor{\mathop{\mathrm{Mor}}\nolimits}
\def\Ob{\mathop{\mathrm{Ob}}\nolimits}
\def\Sh{\mathop{\mathit{Sh}}\nolimits}
\def\NL{\mathop{N\!L}\nolimits}
\def\CH{\mathop{\mathrm{CH}}\nolimits}
\def\proetale{{pro\text{-}\acute{e}tale}}
\def\etale{{\acute{e}tale}}
\def\QCoh{\mathit{QCoh}}
\def\Ker{\mathop{\mathrm{Ker}}}
\def\Im{\mathop{\mathrm{Im}}}
\def\Coker{\mathop{\mathrm{Coker}}}
\def\Coim{\mathop{\mathrm{Coim}}}

% Boxtimes
%
\DeclareMathSymbol{\boxtimes}{\mathbin}{AMSa}{"02}

%
% Macros for moduli stacks/spaces
%
\def\QCohstack{\mathcal{QC}\!\mathit{oh}}
\def\Cohstack{\mathcal{C}\!\mathit{oh}}
\def\Spacesstack{\mathcal{S}\!\mathit{paces}}
\def\Quotfunctor{\mathrm{Quot}}
\def\Hilbfunctor{\mathrm{Hilb}}
\def\Curvesstack{\mathcal{C}\!\mathit{urves}}
\def\Polarizedstack{\mathcal{P}\!\mathit{olarized}}
\def\Complexesstack{\mathcal{C}\!\mathit{omplexes}}
% \Pic is the operator that assigns to X its picard group, usage \Pic(X)
% \Picardstack_{X/B} denotes the Picard stack of X over B
% \Picardfunctor_{X/B} denotes the Picard functor of X over B
\def\Pic{\mathop{\mathrm{Pic}}\nolimits}
\def\Picardstack{\mathcal{P}\!\mathit{ic}}
\def\Picardfunctor{\mathrm{Pic}}
\def\Deformationcategory{\mathcal{D}\!\mathit{ef}}

% --- appended 2026-10-12 (Chief): the project own definitions of the reference/slogan/history
% display environments, copied mechanically from the shared Stacks preamble lines 8-12
% (\usepackage{verbatim} plus the three \newenvironment lines, which the chapter preamble keeps
% verbatim from the project preamble). Nothing added, removed or redefined: these environments
% carry no mathematical content and the project itself typesets them as comments. Without these
% three lines a counted statement whose own text uses them fails with Environment slogan undefined.
\usepackage{verbatim}
\newenvironment{reference}{\comment}{\endcomment}
\newenvironment{slogan}{\comment}{\endcomment}
\newenvironment{history}{\comment}{\endcomment}
\usepackage{cleveref}
\crefname{originallemma}{original lemma}{original lemmas}
\Crefname{originallemma}{Original lemma}{Original lemmas}
\makeatletter
\newlabel{more-algebra-lemma-flat-finite-type-valuation-ring-finite-presentation}{{053E}{source}{}{sourcepointer-1}{}}
\newlabel{more-algebra-lemma-flat-finite-type-valuation-ring-finite-presentation@cref}{{[originallemma][0][]053E}{[1][0][]source}{}{}{}}
\makeatother
\crefname{originallemma}{original lemma}{original lemmas}
\Crefname{originallemma}{Original lemma}{Original lemmas}
\makeatletter
\newlabel{lemma-valuation-ring-torsion-free-flat}{{0539}{source}{}{sourcepointer-2}{}}
\newlabel{lemma-valuation-ring-torsion-free-flat@cref}{{[originallemma][0][]0539}{[1][0][]source}{}{}{}}
\makeatother
\crefname{originallemma}{original lemma}{original lemmas}
\Crefname{originallemma}{Original lemma}{Original lemmas}
\makeatletter
\newlabel{algebra-lemma-homogenize}{{052N}{source}{}{sourcepointer-3}{}}
\newlabel{algebra-lemma-homogenize@cref}{{[originallemma][0][]052N}{[1][0][]source}{}{}{}}
\makeatother
\crefname{originallemma}{original lemma}{original lemmas}
\Crefname{originallemma}{Original lemma}{Original lemmas}
\makeatletter
\newlabel{lemma-torsion}{{0537}{source}{}{sourcepointer-4}{}}
\newlabel{lemma-torsion@cref}{{[originallemma][0][]0537}{[1][0][]source}{}{}{}}
\makeatother
\crefname{originallemma}{original lemma}{original lemmas}
\Crefname{originallemma}{Original lemma}{Original lemmas}
\makeatletter
\newlabel{lemma-flat-graded-finite-type-finite-presentation}{{053D}{source}{}{sourcepointer-5}{}}
\newlabel{lemma-flat-graded-finite-type-finite-presentation@cref}{{[originallemma][0][]053D}{[1][0][]source}{}{}{}}
\makeatother
\crefname{originallemma}{original lemma}{original lemmas}
\Crefname{originallemma}{Original lemma}{Original lemmas}
\makeatletter
\newlabel{lemma-flat-graded-finite-type-finite-presentation-module}{{053C}{source}{}{sourcepointer-6}{}}
\newlabel{lemma-flat-graded-finite-type-finite-presentation-module@cref}{{[originallemma][0][]053C}{[1][0][]source}{}{}{}}
\makeatother
\crefname{originalreference}{original reference}{original references}
\Crefname{originalreference}{Original reference}{Original references}
\makeatletter
\newlabel{Nagata-Finitely}{{cite-Nagata-Finitely}{source}{}{sourcepointer-7}{}}
\newlabel{Nagata-Finitely@cref}{{[originalreference][0][]cite-Nagata-Finitely}{[1][0][]source}{}{}{}}
\makeatother
\title{Flat finite type algebras and flat finite modules over a valuation ring are finitely presented}
\author{The Stacks Project Authors}
\date{stacks-project:more-algebra.tex:local-snapshot-20261012; GNU FDL 1.2 or later}
\begin{document}
\maketitle
\setcounter{subsection}{1}
\noindent\textit{Editorial note.} All delivered mathematical text is the author's own, copied from the pinned Stacks Project chapter file more-algebra.tex; no statement or proof step was rewritten, reordered or removed.\par\medskip
\noindent\textit{Editorial note.} The chapter file pinned-original-source.tex carries only \textbackslash{}input\{preamble\} on its first line, so no preamble text lives inside the chapter. The retained file sources/stacks-more-algebra/retained-preamble.tex is a mechanical, line-for-line copy of the project's single shared preamble as staged in the snapshot (Xy-pic setup plus theorem declarations and mathematical macros). Nothing was added, removed or redefined, and the excerpt serves only as preamble, never as a span coordinate.\par\medskip
\noindent\textit{Editorial note.} Both counted spans start at the \textbackslash{}begin line of their environment, so the statement and the author's proof are delivered whole; surrounding prose of the section is deliberately outside the counted spans and is carried as standard prerequisites instead.\par\medskip
\noindent\textit{Editorial note.} The citation fields (source\_id, source\_version, authors, source\_url, license.name, source\_title) are identical to those of every other card this lane has produced from this chapter, so the batch selector treats them as one source.\par\medskip
\noindent\textit{Allowed standard prerequisites (editorial).}\begin{itemize}
\item Flatness over a valuation ring is the same as torsion freeness. Lemma \textbackslash{}ref\{lemma-valuation-ring-torsion-free-flat\}, cited at proof lines 5981-5982. Lets the author remove torsion parts, which is the key reduction.
\item Homogenisation: a finite type algebra over a ring is the degree zero part of a graded algebra generated in degree one, and similarly for modules. Algebra, Lemma \textbackslash{}ref\{algebra-lemma-homogenize\}, cited at proof lines 5983-5988. Realises the algebra and the module as localisations of graded objects.
\item The torsion submodule can be removed without changing the degree zero part of a graded module. Lemma \textbackslash{}ref\{lemma-torsion\}, cited at proof lines 5989-5992. Makes both the graded algebra and the graded module torsion free.
\item A flat graded algebra of finite type over a local ring and a flat finite graded module over it are finitely presented. Lemma \textbackslash{}ref\{lemma-flat-graded-finite-type-finite-presentation\} and Lemma \textbackslash{}ref\{lemma-flat-graded-finite-type-finite-presentation-module\}, cited at proof lines 5994-5996 and 5999-6001. Give the two graded statements which specialise to the required results.
\end{itemize}\medskip
\ifdefined\hypertarget\hypertarget{sourcepointer-1}{}\fi
\noindent\textit{Source reference.} Original lemma 053E: focused-seeds/AGP/sources/stacks-snapshot-20261011/tags.txt, source lines 6602--6602. The lemma's own canonical tag, quoted so that readers can resolve the label inside the delivered text; it points at the counted statement itself and carries no proof content.\par\medskip
\ifdefined\hypertarget\hypertarget{sourcepointer-2}{}\fi
\noindent\textit{Source reference.} Original lemma 0539: focused-seeds/AGP/sources/stacks-snapshot-20261011/tags.txt, source lines 6597--6597. Navigation pointer to a statement the counted proof refers to but does not contain; nothing load-bearing is delegated to it.\par\medskip
\ifdefined\hypertarget\hypertarget{sourcepointer-3}{}\fi
\noindent\textit{Source reference.} Original lemma 052N: focused-seeds/AGP/sources/stacks-snapshot-20261011/tags.txt, source lines 6576--6576. Navigation pointer to a statement the counted proof refers to but does not contain; nothing load-bearing is delegated to it.\par\medskip
\ifdefined\hypertarget\hypertarget{sourcepointer-4}{}\fi
\noindent\textit{Source reference.} Original lemma 0537: focused-seeds/AGP/sources/stacks-snapshot-20261011/tags.txt, source lines 6595--6595. Navigation pointer to a statement the counted proof refers to but does not contain; nothing load-bearing is delegated to it.\par\medskip
\ifdefined\hypertarget\hypertarget{sourcepointer-5}{}\fi
\noindent\textit{Source reference.} Original lemma 053D: focused-seeds/AGP/sources/stacks-snapshot-20261011/tags.txt, source lines 6601--6601. Navigation pointer to a statement the counted proof refers to but does not contain; nothing load-bearing is delegated to it.\par\medskip
\ifdefined\hypertarget\hypertarget{sourcepointer-6}{}\fi
\noindent\textit{Source reference.} Original lemma 053C: focused-seeds/AGP/sources/stacks-snapshot-20261011/tags.txt, source lines 6600--6600. Navigation pointer to a statement the counted proof refers to but does not contain; nothing load-bearing is delegated to it.\par\medskip
\ifdefined\hypertarget\hypertarget{sourcepointer-7}{}\fi
\noindent\textit{Source reference.} Original reference cite-Nagata-Finitely: sources/stacks-more-algebra/pinned-original-source.tex, source lines 5966--5966. Bibliographic pointer appearing in the counted text; it carries no mathematical argument.\par\medskip
\begin{lemma}
\label{lemma-flat-finite-type-valuation-ring-finite-presentation}
\begin{reference}
\cite[Theorem 3]{Nagata-Finitely}
\end{reference}
Let $A$ be a valuation ring. Let $A \to B$ be a ring map of finite type.
Let $M$ be a finite $B$-module.
\begin{enumerate}
\item If $B$ is flat over $A$, then $B$ is a finitely presented $A$-algebra.
\item If $M$ is flat as an $A$-module, then $M$ is finitely presented
as a $B$-module.
\end{enumerate}
\end{lemma}

\begin{proof}
We are going to use that an $A$-module is flat if and only if it is
torsion free, see
Lemma \ref{lemma-valuation-ring-torsion-free-flat}.
By
Algebra, Lemma \ref{algebra-lemma-homogenize}
we can find a graded $A$-algebra $S$ with $S_0 = A$ and generated
by finitely many elements in degree $1$, an element $f \in S_1$ and a
finite graded $S$-module $N$ such that $B \cong S_{(f)}$ and
$M \cong N_{(f)}$. If $M$ is torsion free, then we can take $N$ torsion
free by replacing it by $N/N_{tors}$, see
Lemma \ref{lemma-torsion}.
Similarly, if $B$ is torsion free, then we can take
$S$ torsion free by replacing it by $S/S_{tors}$.
Hence in case (1), we may apply
Lemma \ref{lemma-flat-graded-finite-type-finite-presentation}
to see that $S$ is a finitely presented
$A$-algebra, which implies that $B = S_{(f)}$ is a finitely
presented $A$-algebra. To see (2) we may first replace $S$ by
a graded polynomial ring, and then we may apply
Lemma \ref{lemma-flat-graded-finite-type-finite-presentation-module}
to conclude.
\end{proof}

\end{document}
